% =============================================================================
\bloco{3}{Integrative: growth, interest rates, money and seigniorage}{Kurlat chs.~9--11; Romer ch.~11 \,\textbullet\, 25 points}

\questao{(10 points) Objective items.}

Items a)--d) use the economy of Question~6 and compare balanced paths, on which $i$, $r$ and $\pi$ are constant. A closed
economy receives the endowment $Y_t=Y_0(1+g)^t$, $g\ge0$. The representative household maximizes
$\sum_{t\ge0}\beta^tc_t^{1-\sigma}/(1-\sigma)$ ($\ln c_t$ if $\sigma=1$), with $\sigma>0$, $0<\beta<1$ and $\beta(1+g)^{1-\sigma}<1$,
and trades one-period real and nominal bonds. Real money demand is $M_t/(P_tY_t)=\kappa e^{-ai_t}$, $\kappa,a>0$; the Fisher
equation holds exactly, $1+i=(1+r)(1+\pi)$; policy keeps $i\ge0$.

\medskip
\textbf{a)} (3 points) \emph{General equilibrium $\times$ money.} Let $\sigma=1$. The central bank adjusts the money supply so
that inflation equals its target $\pi^*>0$. The equilibrium nominal interest rate is:
\begin{alternativas}
\item $i=(1+g)(1+\pi^*)/\beta-1$
\item $i=\beta(1+\pi^*)-1$
\item $i=(1+\pi^*)/\beta-1$
\item $i=\beta(1+g)(1+\pi^*)-1$
\end{alternativas}
\gabarito{A}
\begin{sol}
\passo{Euler equation and goods clearing}{The real-bond first-order condition, $c_t^{-\sigma}=\beta(1+r)c_{t+1}^{-\sigma}$, equates the
marginal utility given up today to the discounted marginal utility of the $1+r$ goods received tomorrow. In a closed endowment
economy no one can lend in the aggregate, so goods clearing forces $c_t=Y_t$ and $c_{t+1}/c_t=1+g$:
\[
1+r=\frac{(1+g)^{\sigma}}{\beta}\ \overset{\sigma=1}{=}\ \frac{1+g}{\beta}.
\]
The real rate is the one at which the household is content to consume exactly the growing endowment.}
\passo{Fisher}{Dividing the nominal-bond condition $c_t^{-\sigma}/P_t=\beta(1+i)c_{t+1}^{-\sigma}/P_{t+1}$ by the real one gives
$1+i=(1+r)(1+\pi)$; with $\pi=\pi^*$, $1+i=(1+g)(1+\pi^*)/\beta$. Money is endogenous here: holding $M/(PY)$ constant at that
$i$ requires $M$ to grow at $(1+g)(1+\pi^*)-1$.}
\passo{distractors}{C uses the no-growth rate $1+r=1/\beta$, ignoring that rising consumption raises $r$; D inverts the discount
factor, writing $\beta u'(c_t)=(1+r)u'(c_{t+1})$, hence $1+r=\beta(1+g)$; B makes both mistakes, so $1+r=\beta$.}
\refbib{Kurlat (2020), ch.~9, \S9.3 (Euler equation); ch.~11, \S11.1 (Fisher equation); Romer (2012), ch.~11, \S11.1.}
\end{sol}

\medskip
\textbf{b)} (3 points) \emph{Seigniorage $\times$ interest rates.} A government finances itself by printing money, as Brazil's
did before the 1994 Real Plan. Let $g=0$ and write $r>0$ for the equilibrium real rate. The money stock grows at a constant
rate $\mu>0$ and seigniorage is $S_t=(M_{t+1}-M_t)/P_t$, so $S/Y=\mu\kappa e^{-ai}$ on a balanced path. The money growth rate
that maximizes $S/Y$ is:
\begin{alternativas}
\item $\mu^*=1/a$
\item $\mu^*=1/[a(1+r)]$
\item $\mu^*=1/a-r$
\item $\mu^*=1/[a(1+r)]-r$
\end{alternativas}
\gabarito{B}
\begin{sol}
\passo{rate and base}{With $g=0$, item a) gives $1+r=1/\beta$, which does not depend on $\mu$. A constant $M/(PY)$ with constant
$Y$ forces $\pi=\mu$, and the exact Fisher equation gives $i=(1+r)(1+\mu)-1$, so $di/d\mu=1+r$.}
\passo{first-order condition}{$\ln(S/Y)=\ln\mu+\ln\kappa-a[(1+r)(1+\mu)-1]$, so
\[
\frac{d\ln(S/Y)}{d\mu}=\frac{1}{\mu}-a(1+r)=0\quad\Longrightarrow\quad\mu^*=\frac{1}{a(1+r)}.
\]}
\passo{second-order condition}{$d^2\ln(S/Y)/d\mu^2=-1/\mu^2<0$ for every $\mu>0$: $\ln(S/Y)$ is strictly concave, so $\mu^*$ is the
unique global maximum; at $\mu^*$, $d^2(S/Y)/d\mu^2=-(S/Y)/\mu^{*2}<0$. Revenue tends to $0$ as $\mu\to0^+$ and as $\mu\to\infty$:
the inflation Laffer curve.}
\passo{distractors}{A uses the approximate Fisher equation $i\approx r+\mu$ ($di/d\mu=1$) and overstates $\mu^*$ by the factor
$1+r$; D subtracts $r$, converting the root as if it were a nominal rate ($\mu=i-r$); C makes both mistakes, subtracting
$r$ from A's root.}
\intuicao{At the peak the base $\kappa e^{-ai}$ has elasticity $-a(1+r)\mu^*=-1$ with respect to $\mu$: a higher tax rate is
exactly offset by the shrinking base. The exact Fisher equation makes the base shrink faster, by the factor $1+r$.}
\refbib{Kurlat (2020), ch.~11, \S11.3 and Exercise~11.6; Romer (2012), ch.~11, \S11.9.}
\end{sol}

\medskip
\textbf{c)} (2 points, T or F) \emph{General equilibrium $\times$ money $\times$ welfare.} In this endowment economy a permanent
rise in money growth leaves the real interest rate unchanged and therefore has no real effects.
\gabarito{F}
\begin{sol}
\passo{what money growth cannot move}{$1+r=(1+g)^\sigma/\beta$ and $c_t=Y_t$ (item a) do not involve $\mu$: consumption and the
real rate are superneutral.}
\passo{what it moves}{$M/(PY)$ is constant on a balanced path, so $1+\pi=(1+\mu)/(1+g)$ and, with $m\equiv\kappa e^{-ai}$,
\[
1+i=\frac{(1+g)^{\sigma-1}(1+\mu)}{\beta},\qquad \frac{\partial i}{\partial\mu}=\frac{(1+g)^{\sigma-1}}{\beta}>0,\qquad
\frac{\partial m}{\partial\mu}=-a\,m\,\frac{\partial i}{\partial\mu}<0 .
\]
Real balances are a real variable, and they fall.}
\passo{welfare}{The Bailey deadweight loss per unit of output, $w(i)$ of Question~6(d), is
\[
w(i)=\int_0^i m(x)\,dx-i\,m(i)=\frac{\kappa}{a}\left[1-(1+ai)e^{-ai}\right],\qquad w'(i)=a\kappa\,i\,e^{-ai}=-i\,m'(i)>0,
\]
and $w(i)>0$ for $i>0$ because $e^{ai}>1+ai$. Higher $\mu$ raises $i$ and the loss; only the Friedman rule $i=0$, that is
$1+\mu=\beta(1+g)^{1-\sigma}$, removes it. Superneutrality fails once holding money yields services.}
\refbib{Kurlat (2020), ch.~10, \S10.4; ch.~11, \S11.2 (neutrality) and \S11.4 (cost of inflation, Friedman rule).}
\end{sol}

\medskip
\textbf{d)} (2 points, T or F) \emph{Growth $\times$ money.} Under a constant money-growth rule, faster trend growth always
lowers the nominal interest rate because it lowers inflation.
\gabarito{F}
\begin{sol}
\passo{two channels}{From c), $1+\pi=(1+\mu)/(1+g)$ falls with $g$ (growing demand for real balances absorbs part of the money
growth), while $1+r=(1+g)^\sigma/\beta$ rises with $g$. From $1+i=(1+g)^{\sigma-1}(1+\mu)/\beta$:
\[
\frac{\partial\ln(1+i)}{\partial\ln(1+g)}=\underbrace{\sigma}_{\text{real rate}}\ \underbrace{-\ 1}_{\text{inflation}}=\sigma-1 .
\]}
\passo{sign}{$i$ rises with $g$ if $\sigma>1$, does not depend on $g$ if $\sigma=1$, and falls only if $\sigma<1$, that is, only
when the elasticity of intertemporal substitution $1/\sigma$ exceeds one. The log case already refutes ``always''.}
\intuicao{Under inflation targeting $\partial\ln(1+i)/\partial\ln(1+g)=\sigma>0$ for every $\sigma$: the money rule passes faster
growth into lower inflation one for one, but with $\sigma>1$ the real rate rises by more.}
\refbib{Kurlat (2020), ch.~9, \S9.3; ch.~11, \S11.2 (steady-state inflation); Romer (2012), ch.~11, \S11.1.}
\end{sol}
